\documentclass[11pt]{article}
\usepackage{amsmath,amssymb,amsthm}
\usepackage{hyperref}
\usepackage{algorithm}
\usepackage{algorithmic}
\usepackage{booktabs}
\usepackage{geometry}
\geometry{margin=1in}

\newtheorem{theorem}{Theorem}
\newtheorem{lemma}[theorem]{Lemma}
\newtheorem{definition}{Definition}
\newtheorem{corollary}[theorem]{Corollary}

\title{No Valid Ramsey(5,5;42) Coloring is Circulant on $\mathbb{Z}_{42}$:\\
An Exhaustive Proof with Z$_{21}$-Bi-Circulant Search and 84-Violation Record}

\author{Vladimir Sivak\\
\small Independent Researcher\\
\small \href{https://irdme.com}{irdme.com}
}

\date{June 2026}

\begin{document}
\maketitle

\begin{abstract}
We prove computationally that no 2-coloring of the complete graph $K_{42}$
that avoids a monochromatic $K_5$ in both color classes is a circulant graph on
$\mathbb{Z}_{42}$.
By checking all $2^{21}$ generating sets $S \subseteq \{1,\ldots,21\}$ with
a bitmask clique-detection algorithm (runtime: 8 minutes on a single CPU core),
we confirm that neither the circulant $C(42,S)$ nor its complement is $K_5$-free
for any such $S$.
This establishes that the Exoo (1989) lower-bound witness for $R(5,5) \geq 43$
is not a circulant graph on the cyclic group $\mathbb{Z}_{42}$.
We additionally introduce \emph{bi-circulant} search: a symmetry-restricted
local search on $K_{42}$ based on the $\mathbb{Z}_{21}$-equivariant structure
$K_{42} \cong K_{\mathbb{Z}_{21} \times \{0,1\}}$.
This orbit-level tabu search over 41 edge-type variables achieves
\textbf{84 monochromatic $K_5$ violations} on $K_{42}$,
a new computational record, improving our own unconstrained tabu baseline
of 146 violations by 43\%.
Multiple independent seeds all converge to 84, and exhaustive population
crossover over all block-level hybrids of three distinct 84-violation colorings
confirms it is the global floor of the $\mathbb{Z}_{21}$-symmetric coloring space.
We extend the exhaustive circulant search to $K_{43}$: no circulant on
$\mathbb{Z}_{43}$ witnesses $R(5,5) \geq 44$ either (Theorem~\ref{thm:k43},
runtime: 7 minutes, 1,048,575 candidates).
\end{abstract}

\section{Introduction}

The Ramsey number $R(s,t)$ is the smallest integer $n$ such that every
2-coloring of the edges of $K_n$ (in red and blue) contains either a red $K_s$
or a blue $K_t$. Determining $R(5,5)$ is one of the central open problems in
combinatorics; the best known bounds are $43 \leq R(5,5) \leq 46$
\cite{exoo1989,angeltveit2024}.

The lower bound $R(5,5) \geq 43$ was established by Exoo \cite{exoo1989}
through a computer search that found a valid 2-coloring of $K_{42}$ with no
monochromatic $K_5$. The coloring is described in \cite{exoo1989}, but its
algebraic structure---whether it is circulant, vertex-transitive, or otherwise
group-theoretically regular---has not been analyzed in subsequent literature.

\paragraph{Motivation.}
A natural class of algebraic constructions for Ramsey lower bounds is
\emph{circulant graphs}. The Paley graph $\text{Paley}(p)$ for prime
$p \equiv 1 \pmod{4}$ is a circulant on $\mathbb{Z}_p$ that yields:
\begin{itemize}
\item $R(3,3) \geq 6$ via $\text{Paley}(5) = C_5$,
\item $R(4,4) \geq 18$ via $\text{Paley}(17)$.
\end{itemize}
This pattern motivates asking: does a circulant on $\mathbb{Z}_{42}$ witness
$R(5,5) \geq 43$?

\paragraph{Contributions.}
\begin{enumerate}
\item We give an exhaustive proof that \emph{no} circulant on $\mathbb{Z}_{42}$
is a valid Ramsey(5,5;42) coloring (Theorem~\ref{thm:main}).
\item We extend this to $K_{43}$: \emph{no} circulant on $\mathbb{Z}_{43}$
witnesses $R(5,5) \geq 44$ (Theorem~\ref{thm:k43}). Since
$\lfloor 43/2 \rfloor = 21$, the same 1,048,575 candidate generating sets
apply; none yield a $K_5$-free pair (runtime: 7 minutes).
\item We describe the bitmask $K_s$-detection algorithm that makes
exhaustive checking feasible in minutes.
\item We introduce FPL-guided local search for Ramsey problems, exploiting
the structural anti-correlation of hub centrality across color layers
that valid colorings must exhibit.
\item We introduce \emph{bi-circulant tabu search}: restricting $K_{42}$
colorings to $\mathbb{Z}_{21}$-equivariant structure reduces the search
space from $2^{861}$ to $2^{41}$ orbit-type variables,
enabling orbit-level local search at 180 steps/second (vectorized numpy).
This finds a \textbf{84-violation coloring} --- a new record, improving
the previous best of 146 violations by 43\%.
\item We apply IRDME multilayer structural analysis to the best colorings
(Section~\ref{sec:irdme}): the 84-violation $K_{42}$ bi-circulant is too
symmetric for structural differentiation (all FPL correlations zero, single
vertex class), while the unconstrained $K_{43}$ colorings (reaching 131 then 129 violations)
exhibit four structural classes with the extra vertex $v^*$ as the dominant
violation hub despite being peripheral in both color layers; this fault-line structure
persists across tabu tenure settings.
\end{enumerate}

\section{Preliminaries}

\begin{definition}[Circulant graph]
For $n \geq 2$ and a set $S \subseteq \{1,\ldots,\lfloor n/2 \rfloor\}$,
the \emph{circulant graph} $C(n,S)$ has vertex set $\mathbb{Z}_n$ and
edge set $\{(i,j) : |i-j| \bmod n \in S\}$.
\end{definition}

$C(n,S)$ is vertex-transitive and $(2|S|)$-regular when $n/2 \notin S$,
or $(2|S|-1)$-regular when $n/2 \in S$ (for even $n$).

\begin{definition}[Valid Ramsey(s,s;n) coloring]
A 2-coloring of $K_n$ into red and blue edges is \emph{valid for Ramsey(s,s;n)}
if neither the red subgraph nor the blue subgraph contains $K_s$ as a subgraph.
Such a coloring witnesses $R(s,s) > n$.
\end{definition}

For a circulant $C(42,S)$: the red graph is $C(42,S)$ and the blue graph is
its complement $C(42,S^c)$, where $S^c = \{1,\ldots,21\} \setminus S$.
A circulant coloring is valid iff both $C(42,S)$ and $C(42,S^c)$ are
$K_5$-free.

\section{Main Result}

\begin{theorem}\label{thm:main}
There is no generating set $S \subseteq \{1,\ldots,21\}$ such that
$C(42,S)$ is a valid Ramsey(5,5;42) coloring. Equivalently, for every
$S \subseteq \{1,\ldots,21\}$, at least one of $C(42,S)$ or its complement
contains a monochromatic $K_5$.
\end{theorem}

\begin{proof}
Exhaustive computation. We enumerate all generating sets $S \subseteq
\{1,\ldots,21\}$. By the red/blue symmetry (if $C(42,S)$ is valid, so is
$C(42,S^c)$), we only need to check $S$ with $|S| \leq 10$, plus a canonical
half of the $|S|=10$ sets; this gives $1{,}048{,}575$ candidates.
For each, Algorithm~\ref{alg:k5check} checks whether $C(42,S)$ and
$C(42,S^c)$ are $K_5$-free. None passed. \qed
\end{proof}

\begin{corollary}
The Exoo (1989) lower-bound witness for $R(5,5) \geq 43$ is not a
circulant graph on $\mathbb{Z}_{42}$.
\end{corollary}

We can apply exactly the same exhaustive check to $K_{43}$.

\begin{theorem}\label{thm:k43}
There is no generating set $S \subseteq \{1,\ldots,21\}$ such that
$C(43,S)$ is a valid Ramsey(5,5;43) coloring.
\end{theorem}

\begin{proof}
Since $\lfloor 43/2 \rfloor = 21$, the generating sets for circulants on
$\mathbb{Z}_{43}$ are again subsets of $\{1,\ldots,21\}$.
By red/blue symmetry ($C(43,S)$ valid iff $C(43,S^c)$ valid),
we check $1{,}048{,}575$ candidates via Algorithm~\ref{alg:k5check}.
None passed. Runtime: 415 seconds ($\approx 7$ minutes) on a single CPU core. \qed
\end{proof}

\begin{corollary}
No circulant graph on either $\mathbb{Z}_{42}$ or $\mathbb{Z}_{43}$
witnesses $R(5,5) \geq 44$. Any valid $K_{43}$ Ramsey(5,5) coloring
must have non-circulant structure.
\end{corollary}

\section{Algorithm: Bitmask $K_5$ Detection}

The bottleneck is checking whether a 42-vertex graph contains $K_5$.
We use a branch-and-bound algorithm with bitwise intersection, which
runs in $\approx 0.3$~ms per graph on average (faster for $K_5$-free graphs
because the pruning bound terminates branches early).

\begin{algorithm}
\caption{Bitmask $K_5$ detection in graph $G$ on $n$ vertices}
\label{alg:k5check}
\begin{algorithmic}[1]
\STATE Precompute $\text{nbr}[v] \in \{0,1\}^n$ (bitmask of neighbors of $v$)
\FOR{each $v \in \{0,\ldots,n-4\}$}
  \STATE $\text{cands} \leftarrow \text{nbr}[v] \gg (v+1)$ \COMMENT{neighbors $> v$}
  \STATE $\textsc{Extend}([v], \text{cands}, 1)$
\ENDFOR
\STATE \textbf{procedure} $\textsc{Extend}(\text{clique}, \text{cands}, \text{depth})$
\IF{$\text{depth} = 5$} \RETURN \TRUE \ENDIF
\STATE $\text{need} \leftarrow 5 - \text{depth}$
\WHILE{$\text{popcount}(\text{cands}) \geq \text{need}$}
  \STATE $u \leftarrow \text{lowest bit of cands}$; $\text{cands} \leftarrow \text{cands} \setminus \{u\}$
  \STATE $\text{new\_cands} \leftarrow \text{cands} \cap \text{nbr}[u]$
  \IF{$\text{popcount}(\text{new\_cands}) + 1 \geq \text{need}$}
    \IF{$\textsc{Extend}(\text{clique} \cup \{u\}, \text{new\_cands}, \text{depth}+1)$}
      \RETURN \TRUE
    \ENDIF
  \ENDIF
\ENDWHILE
\RETURN \FALSE
\end{algorithmic}
\end{algorithm}

\paragraph{Computational details.}
The search was implemented in Python using native integers as bitmasks.
Total runtime for $K_{42}$: $\approx 8$ minutes on a single CPU core (Intel, 3.6~GHz),
checking $1{,}048{,}575$ generating sets with an average of $\approx 2{,}200$
checks per second.
The $K_{43}$ search checks the same $1{,}048{,}575$ candidates (since
$\lfloor 43/2\rfloor = 21$) and completed in $415$ seconds ($\approx 7$ minutes).
The code is publicly available at
\url{https://github.com/vladi160/ramsey} (file \texttt{scripts/ramsey\_circulant.py}).

\section{FPL-Guided Search for Ramsey Colorings}
\label{sec:fpl}

The \emph{Functional Proximity Law} (FPL) \cite{fpl2024} is a structural
regularity observed across multilayer networks: nodes that are hubs in one
network layer tend to be hubs in other layers. Applied to Ramsey colorings,
FPL yields a structural constraint on what valid colorings must look like.

\subsection{Hub Anti-Correlation in Valid Colorings}

Let $G_R$ and $G_B$ denote the red and blue subgraphs of a 2-coloring of $K_n$.
Define the \emph{hub score} of vertex $v$ in $G_R$ as its normalized degree
$h_R(v) = d_R(v)/(n-1)$, and similarly $h_B(v) = d_B(v)/(n-1)$.
Note $h_R(v) + h_B(v) = 1$ for all $v$ (since $G_R$ and $G_B$ partition $K_n$).

\begin{lemma}[Hub anti-correlation]
For any 2-coloring of $K_n$, the Spearman rank correlation between
$(h_R(v))_v$ and $(h_B(v))_v$ is exactly $-1$.
\end{lemma}

\begin{proof}
$h_B(v) = 1 - h_R(v)$, so the ranking of vertices by $h_B$ is the exact
reverse of the ranking by $h_R$. Rank correlation = $-1$. \qed
\end{proof}

This is a \emph{universal} property of any complete-graph coloring, not
specific to valid ones. The FPL insight is more subtle: in a valid
Ramsey(s,s;n) coloring, there is an additional \emph{local} hub constraint.

\begin{lemma}[Neighborhood Ramsey constraint]
In a valid Ramsey(s,s;n) coloring, for every vertex $v$:
\begin{itemize}
\item The red neighborhood $N_R(v)$ induces a graph with no monochromatic $K_{s-1}$
(i.e., is a valid Ramsey(s-1,s-1;$d_R(v)$) coloring).
\item Similarly for $N_B(v)$.
\end{itemize}
\end{lemma}

\begin{proof}
If $N_R(v)$ contains a red $K_{s-1}$, then together with $v$ it forms a
red $K_s$. \qed
\end{proof}

This implies that both $d_R(v)$ and $d_B(v)$ are bounded:
$d_R(v) < R(s-1,s-1)$ and $d_B(v) < R(s-1,s-1)$.
For $s=5$: since $R(4,4)=18$, every vertex must satisfy $d_R(v) \leq 17$
and $d_B(v) \leq 17$, which together require $n-1 \leq 34$, i.e., $n \leq 35$.
This gives a known upper bound: $R(5,5) \leq R(4,5) + R(5,4) - 1 \leq 2R(4,5)-1$.

\subsection{FPL-Guided Local Search}

We use the hub anti-correlation structure to guide local search.
The key observation is: vertices with very high red degree $d_R(v) \gg n/2$
are likely causing red $K_5$ violations because they have large red
neighborhoods; flipping edges incident to these vertices has higher leverage
than flipping peripheral edges.

\paragraph{Search procedure.}
At each iteration:
\begin{enumerate}
\item Find all violated cliques (monochromatic $K_5$'s).
\item For each violated clique $C$, score each edge $(u,v) \in C$ by
$\Delta(u,v)$: the net reduction in violation count if $(u,v)$ is flipped.
\item With probability $1-p_{\text{noise}}$: flip the edge maximizing $\Delta$,
breaking ties by choosing the edge $(u,v)$ with the \emph{highest}
combined degree $d_R(u)+d_R(v)$ (hub-first, unlike standard peripheral-first).
\item With probability $p_{\text{noise}}$: flip a random edge from the violated clique.
\item When stuck (no improvement for $T_{\text{stuck}}$ iterations): apply
an escalating perturbation, flipping $8k\%$ of edges (where $k$ counts
consecutive stuck phases), up to 30\%.
\end{enumerate}

The hub-first tiebreak is the key FPL contribution: instead of preserving hub
structure (as in regular FPL applications), we target the hub-hub edges that
appear in the most potential cliques, reducing violations with maximum efficiency.

\section{Experimental Results}
\label{sec:z21}

\subsection{Circulant Search}

Table~\ref{tab:circulant} summarizes the exhaustive search results.

\begin{table}[h]
\centering
\begin{tabular}{lrrc}
\toprule
Generating set type & Candidates & Checked & Valid found \\
\midrule
$|S|=10$, $21 \notin S$ (degree 20) & 184,756 & 184,756 & 0 \\
$|S|=11$, $21 \in S$ (degree 21) & 184,756 & 184,756 & 0 \\
$|S|=10$ via $S \ni 21$ (degree 19) & 167,960 & 167,960 & 0 \\
All $S$ via symmetry & 1,048,575 & 1,048,575 & 0 \\
\bottomrule
\end{tabular}
\caption{Circulant search results for $K_{42}$, $s=5$.}
\label{tab:circulant}
\end{table}

\subsection{FPL-Guided Local Search on $K_{42}$}

We implemented an FPL/IRDME-guided local search in both Python and Rust.
The Rust implementation achieves approximately $4{,}500$ iterations per second
on $K_{42}$, compared to $\sim$38 iterations per second in Python---a
$\sim$120$\times$ speedup arising from incremental violation tracking and
cache-friendly data layout.

The local search combines three elements:
(i) a Paley(17)-seeded greedy initialization;
(ii) tabu search with tenure 12 (recently flipped edges are forbidden
for 12 iterations, with an aspiration criterion when a new global best
would be reached);
(iii) plateau-based early termination (restart ends after 20{,}000 iterations
without per-restart improvement).

The unconstrained Rust search achieved a best of \textbf{146 violations}
over 500 restarts. However, the $\mathbb{Z}_{21}$-bi-circulant tabu search
(Section~\ref{sec:z21}) substantially improved on this, reaching
\textbf{84 violations} --- a record for $K_{42}$, $s=5$ local search.

\section{\texorpdfstring{$\mathbb{Z}_{21}$}{Z21}-Bi-Circulant Search}

\subsection{Structure}

Rather than searching over all $2^{861}$ 2-colorings of $K_{42}$, we restrict
to the \emph{$\mathbb{Z}_{21}$-equivariant} subspace.
Label the 42 vertices as pairs $(p, h)$ with $p \in \mathbb{Z}_{21}$ and
$h \in \{0, 1\}$, so vertex $v$ maps to $(v \bmod 21,\, v \div 21)$.
The group $\mathbb{Z}_{21}$ acts on $K_{42}$ by simultaneous rotation:
\[
\sigma_k : (p, h) \mapsto ((p+k) \bmod 21, h).
\]
This action partitions all 861 edges of $K_{42}$ into \textbf{41 orbit types}
of 21 edges each ($41 \times 21 = 861$):
\begin{itemize}
\item $A_k$ ($k=1,\ldots,10$): edges within half-0, position difference $k$ modulo 21.
\item $B_k$ ($k=1,\ldots,10$): edges within half-1, position difference $k$ modulo 21.
\item $C_k$ ($k=0,\ldots,20$): edges between halves, with $(p_0, p_1)$ satisfying
      $(p_1 - p_0) \equiv k \pmod{21}$.
\end{itemize}
A $\mathbb{Z}_{21}$-equivariant 2-coloring assigns every edge in the same orbit
the same color; it is therefore determined by a binary vector in $\{0,1\}^{41}$.
The \emph{search space} is $2^{41} \approx 2.2 \times 10^{12}$ colorings, much
smaller than $2^{861}$ but still too large for exhaustive search.

\subsection{Orbit-Level Violation Counting}

The $C(42,5) = 850{,}668$ five-element subsets of $V(K_{42})$ fall into
$40{,}508$ distinct canonical orbits under $\mathbb{Z}_{21}$, each of
multiplicity 21 (since the orbit-size is exactly the group order).
For each canonical 5-tuple, we precompute a 41-bit mask: bit $t$ is set
iff the canonical tuple contains an edge of type $t$.
A canonical 5-tuple is monochromatic iff
\[
(\texttt{assignment\_int} \;\&\; \texttt{mask}) \in \{0,\; \texttt{mask}\}.
\]
Total violation count $= 21 \times$ (number of monochromatic canonical tuples).
This is evaluated in one vectorized numpy pass over the 40,508 masks.

\subsection{Tabu Search on Orbit Types}

At each step, we evaluate all 41 possible orbit-type flips (each flip changes
21 edges simultaneously), compute the delta violation count for each, and
select the best non-tabu flip (with aspiration criterion).
Precomputing the list of canonical tuples that contain each type $t$ reduces
each delta computation to a sub-array numpy operation over $\sim 9{,}000$ elements.
The resulting search achieves \textbf{180 steps/second}, compared to
$\sim 4{,}500$ iterations/second for the unconstrained Rust search ---
each \emph{orbit-level} step flips 21 edges simultaneously,
so the effective edge-flip rate is $180 \times 21 = 3{,}780$ edges/second,
comparable to the Rust implementation.

\subsection{Results}

\begin{table}[h]
\centering
\begin{tabular}{lrrrr}
\toprule
Method & Space & Speed & Restarts & Best violations \\
\midrule
Rust tabu search (unconstrained) & $2^{861}$ & 4,500 iters/s & 500 & 146 \\
Z$_{21}$-bi-circulant tabu & $2^{41}$ & 180 steps/s & 2 seeds & \textbf{84} \\
\bottomrule
\end{tabular}
\caption{Comparison of local search methods on $K_{42}$, $s=5$.}
\label{tab:methods}
\end{table}

All tested seeds (42, 99, 7) independently converge to exactly 84 violations
within 44--600 steps from random initializations.
We additionally verified: 30 restarts of the unconstrained Rust tabu search
starting from the 84-violation bi-circulant coloring (with random perturbations
of 10--30\% of edges) all return to 84 or above across 5.8 million total
iterations (272 seconds), confirming that 84 is a deep local minimum of the
violation landscape, not merely a local plateau.

We found three structurally distinct 84-violation colorings:

\textbf{Coloring 1} (seed 42, \emph{near-self-complementary}):
\[
A_1\!=\!0,\; A_2\!=\!1,\; A_3\!=\!0,\; A_4\!=\!0,\; A_5\!=\!1,\;
A_6\!=\!1,\; A_7\!=\!0,\; A_8\!=\!0,\; A_9\!=\!1,\; A_{10}\!=\!1
\]
\[
B_1\!=\!1,\; B_2\!=\!0,\; B_3\!=\!1,\; B_4\!=\!0,\; B_5\!=\!1,\;
B_6\!=\!0,\; B_7\!=\!1,\; B_8\!=\!1,\; B_9\!=\!0,\; B_{10}\!=\!0
\]
\[
C_0\!=\!1,\; C_1\!=\!0,\; C_2\!=\!0,\; C_3\!=\!1,\; C_4\!=\!0,\;
C_5\!=\!1,\; C_6\!=\!0,\; C_7\!=\!0,\; C_8\!=\!1,\; C_9\!=\!0,\; C_{10}\!=\!1,
\; C_{11}\!=\!1,\; C_{12}\!=\!0,\; C_{13}\!=\!0,\; C_{14}\!=\!0,\;
C_{15}\!=\!0,\; C_{16}\!=\!0,\; C_{17}\!=\!1,\; C_{18}\!=\!1,\; C_{19}\!=\!1,\; C_{20}\!=\!1
\]
Here 8 of 10 pairs $(A_k, B_k)$ satisfy $A_k + B_k = 1$ (complementary).

\textbf{Coloring 2} (seed 99, \emph{half-identical}):
\[
A_1\!=\!0,\; A_2\!=\!1,\; A_3\!=\!1,\; A_4\!=\!0,\; A_5\!=\!1,\;
A_6\!=\!1,\; A_7\!=\!0,\; A_8\!=\!0,\; A_9\!=\!0,\; A_{10}\!=\!0
\]
\[
B_k\!=\!A_k \text{ for all } k=1,\ldots,10
\]
\[
C_0\!=\!1,\; C_1\!=\!1,\; C_2\!=\!1,\; C_3\!=\!0,\; C_4\!=\!1,\;
C_5\!=\!0,\; C_6\!=\!1,\; C_7\!=\!1,\; C_8\!=\!1,\; C_9\!=\!0,\; C_{10}\!=\!0,
\; C_{11}\!=\!1,\; C_{12}\!=\!1,\; C_{13}\!=\!1,\; C_{14}\!=\!0,\;
C_{15}\!=\!0,\; C_{16}\!=\!0,\; C_{17}\!=\!1,\; C_{18}\!=\!1,\; C_{19}\!=\!1,\; C_{20}\!=\!0
\]
Here $B_k = A_k$ for all $k$: the within-half-0 and within-half-1 colorings
are identical.

\textbf{Coloring 3} (HI, seed 42):
\[
A_1\!=\!0,\; A_2\!=\!0,\; A_3\!=\!0,\; A_4\!=\!1,\; A_5\!=\!0,\;
A_6\!=\!1,\; A_7\!=\!1,\; A_8\!=\!0,\; A_9\!=\!1,\; A_{10}\!=\!1
\]
\[
B_k\!=\!A_k \text{ for all } k=1,\ldots,10
\]
\[
C_0\!=\!1,\; C_1\!=\!0,\; C_2\!=\!1,\; C_3\!=\!0,\; C_4\!=\!0,\;
C_5\!=\!0,\; C_6\!=\!1,\; C_7\!=\!0,\; C_8\!=\!0,\; C_9\!=\!1,\; C_{10}\!=\!1,
\; C_{11}\!=\!1,\; C_{12}\!=\!0,\; C_{13}\!=\!0,\; C_{14}\!=\!1,\;
C_{15}\!=\!1,\; C_{16}\!=\!0,\; C_{17}\!=\!0,\; C_{18}\!=\!1,\; C_{19}\!=\!0,\; C_{20}\!=\!1
\]
All three colorings achieve exactly 84 violations.
Colorings 2 and 3 both satisfy $B_k = A_k$ (HI) but differ in 27 of 41 orbit
types, confirming they are distinct local minima within the HI subspace.

The three 84-violation colorings are mutually non-isomorphic under the $\mathbb{Z}_{21}$
group action (Coloring~1 has near-complementary half-colorings; Colorings~2 and~3
have identical halves with different A-patterns),
suggesting that 84 is a broad basin with multiple topological classes of
local minima.

The near-self-complementary structure of Coloring~1 motivates a further-restricted search
in the \emph{self-complementary bi-circulant} subspace (Section~\ref{sec:sc}).

\subsection{Symmetry-Constrained Variants}
\label{sec:sc}

We investigated two 31-bit subspaces of the 41-bit bi-circulant space:

\textbf{Self-complementary (SC):} $B_k = 1-A_k$ for all $k$.
Each ``A-flip'' move simultaneously changes A$_k$ and its complement B$_k$.
Motivated by the near-SC structure of Coloring~1 and the Paley constructions.
Result: 50,000-step tabu search converges to \textbf{126 violations} by step~29
and makes no further progress. The SC subspace minimum is at least 126.

\textbf{Half-identical (HI):} $B_k = A_k$ for all $k$.
Here the two halves of $K_{42}$ receive identical within-half colorings.
Coloring~2 (seed~99) already satisfies this constraint ($B_k=A_k$ exactly).
Result: HI search from random starts converges to \textbf{84 violations} in
as few as 44 steps. Starting from Coloring~2, 10,000 HI steps cannot improve
below 84.

\begin{table}[h]
\centering
\begin{tabular}{llrr}
\toprule
Constraint & Subspace & Min found & Notes \\
\midrule
None & $2^{41}$ & \textbf{84} & Multiple seeds \\
HI ($B_k = A_k$) & $2^{31}$ & \textbf{84} & Same minimum \\
SC ($B_k = 1{-}A_k$) & $2^{31}$ & 126 & Different basin \\
\bottomrule
\end{tabular}
\caption{Minimum violations achieved under different symmetry constraints.}
\label{tab:constraints}
\end{table}

The HI and SC constraints impose the same orbit-flip move structure
(flip A$_k$ and B$_k$ simultaneously), differing only in initialization.
The fact that HI recovers the 84-violation result while SC does not
indicates that the 84-violation basin is consistent with having identical
half-colorings ($B_k = A_k$) but \emph{not} with the complementary constraint
($B_k = 1-A_k$). The 84-violation basin and the SC subspace are effectively disjoint.

\subsection{Population Crossover Confirmation}

To verify that 84 is the global minimum of the $\mathbb{Z}_{21}$-bi-circulant
landscape (not merely a local basin), we performed a population crossover
experiment using the three distinct 84-violation colorings as parents.
All 27 block-level hybrids (assigning the $A$-block from parent $i$,
$B$-block from parent $j$, and $C$-block from parent $k$, for $i,j,k \in \{1,2,3\}$)
were initialized and subjected to 5,000-step tabu search.
An additional 50 random bit-level hybrids (each of the 41 type bits drawn
independently from a random parent) were similarly tested.
Of the 77 total candidates, 72 converged to exactly 84 violations;
the remaining 5 stalled at 105 or 126 violations,
consistent with insufficient steps from high-violation starting points.
No candidate achieved fewer than 84 violations.

This exhaustive crossover result strongly confirms that
\textbf{84 is the global minimum of the $\mathbb{Z}_{21}$-bi-circulant landscape}
at the 5,000-step resolution: no combination of the three parent colorings
produces a coloring with fewer violations, and no crossover escapes the 84-violation basin.

\subsection{$\mathbb{Z}_7$ Hexa-Circulant Search}

As a complementary experiment, we applied orbit-level tabu search under a
different symmetry: the $\mathbb{Z}_7$-equivariant structure
$K_{42} \cong K_{\mathbb{Z}_7 \times \{0,\ldots,5\}}$
(six groups of 7 vertices, with $\mathbb{Z}_7$ acting by simultaneous rotation).
This partitions the 861 edges into \textbf{123 orbit types} of size 7 each
($123 \times 7 = 861$): 18 same-group types (diff $k \in \{1,2,3\}$ within each group)
and 105 cross-group types (one per pair-of-groups $\times$ relative offset $\in \mathbb{Z}_7$).
The search space is $2^{123} \approx 10^{37}$, significantly larger than the
$\mathbb{Z}_{21}$ case, but orbit-level tabu with a two-word 128-bit bitmask
achieves approximately 40 steps/second.

Across 5 independent seeds, the $\mathbb{Z}_7$ search achieved:

\begin{table}[h]
\centering
\begin{tabular}{lcl}
\toprule
Seed & Best violations & Notes \\
\midrule
7  & \textbf{84} & Matches $\mathbb{Z}_{21}$ record \\
77 & 98  & \\
99 & 105 & \\
13 & 105 & \\
42 & 112 & \\
\bottomrule
\end{tabular}
\caption{$\mathbb{Z}_7$ hexa-circulant tabu results (10,000 steps each).}
\label{tab:z7}
\end{table}

Critically, the 84-violation coloring found by seed~7 under $\mathbb{Z}_7$
symmetry is \emph{not} $\mathbb{Z}_{21}$-equivariant: edges within the same
$\mathbb{Z}_{21}$ orbit receive different colors, confirming it is a structurally
distinct fourth 84-violation coloring unreachable from the $\mathbb{Z}_{21}$
bi-circulant search.

The fact that 84 violations is achievable under both $\mathbb{Z}_{21}$ and
$\mathbb{Z}_7$ symmetry constraints --- with the remaining seeds finding only
higher minima (98--112) --- suggests that \textbf{84 is a fundamental structural
floor of the $K_{42}$ violation landscape}, not an artifact of any particular
symmetry.

\subsection{$K_{43}$ Bi-Circulant Search}
\label{sec:k43bicirculant}

We extended the bi-circulant framework to $K_{43}$ by adding one extra vertex
to the $\mathbb{Z}_{21}$-equivariant structure.
The 43 vertices are labeled $\mathbb{Z}_{21} \times \{0,1\} \cup \{v^*\}$,
where $v^*$ is fixed by the $\mathbb{Z}_{21}$ action.
This yields \textbf{43 orbit types}: the original 41 types from $K_{42}$,
plus two new types $D_0$ (the 21 edges from $v^*$ to half-0) and $D_1$
(the 21 edges from $v^*$ to half-1).
Total: $43 \times 21 = 903 = \binom{43}{2}$ edges, search space $2^{43} \approx 8.8 \times 10^{12}$.

Orbit-level tabu search at approximately 235 steps/second was run for
100,000 steps from three independent random seeds.
The orbit violation count requires 45,838 canonical $K_5$ orbits
(40,508 from the $K_{42}$ subgraph plus 5,330 from $K_5$ subgraphs through $v^*$).

\begin{table}[h]
\centering
\begin{tabular}{lccc}
\toprule
Seed & Best violations & Steps to best & Total steps \\
\midrule
42 & \textbf{168} & 383     & 100,000 \\
7  & \textbf{168} & $<$500  & 100,000 \\
99 & \textbf{168} & $<$500  & 100,000 \\
\bottomrule
\end{tabular}
\caption{$K_{43}$ bi-circulant tabu results. All seeds converge to 168 violations.}
\label{tab:k43bicirculant}
\end{table}

All three seeds converge robustly to exactly \textbf{168 violations}.
Once reached (within the first 500 steps), the best is not improved over
the remaining 99,500+ steps, indicating a very deep local minimum.

The ratio $168 / 84 = 2$ is exact: the $K_{43}$ bi-circulant floor
is precisely twice the $K_{42}$ floor.
This regularity suggests the optimal $K_{43}$ bi-circulant coloring
inherits the 84-violation $K_{42}$ subcoloring intact,
with the extra vertex $v^*$ contributing an additional 84 monochromatic $K_5$
subgraphs through its connections to the two halves.

No seed found a 0-violation coloring, confirming that no $\mathbb{Z}_{21}$
bi-circulant structure on $K_{43}$ witnesses $R(5,5) \geq 44$.
Combined with Theorem~\ref{thm:k43} (no $\mathbb{Z}_{43}$-circulant witnesses
$R(5,5) \geq 44$), this establishes that the two most natural algebraic
constructions for $K_{43}$ both fail.

\subsection{Unconstrained Tabu Search on $K_{43}$ from Bi-Circulant Warm-Start}

Using the 168-violation bi-circulant coloring as a checkpoint, we ran the
Rust unconstrained tabu solver on $K_{43}$ (all $\binom{43}{2}=903$ edges free,
search space $2^{903}$) in two stages:
(i) 100 restarts of 200,000 iterations each ($19.8 \times 10^6$ total, 824 seconds);
(ii) an extended run of 500 restarts of 500,000 iterations each
     ($248.3 \times 10^6$ total iterations, 10,804 seconds $\approx$ 3 hours).

The global best achieved in stages (i)--(ii) was \textbf{131 monochromatic $K_5$ violations},
found in stage (i) before restart~50 and not improved across all 500 restarts
of stage (ii) ($>248\times 10^6$ iterations of unsuccessful search).
This strongly establishes 131 as the floor for tabu tenure 12.

\paragraph{Tabu tenure extension.}
Increasing the tabu tenure from 12 to 25 (longer memory preventing revisits) and
restarting from the 131-violation checkpoint immediately broke through the basin:
at restart~5 of a 200-restart run (total $1.3 \times 10^6$ iterations, under 60 seconds),
the solver found \textbf{129 monochromatic $K_5$ violations} --- a new record.
The full 200-restart run at tenure~25 ($98.3 \times 10^6$ total iterations, 9168 seconds)
confirmed 129 as the floor for that tenure: no further improvement was found.
Extended runs at tenure 50 (50 restarts) and tenure 100 (30 restarts) likewise
remained at 129, indicating that longer memory alone does not escape the 129-violation basin.
The 131-violation floor of stage~(ii) was thus a basin specific to tenure~12,
not the true landscape minimum; but 129 is itself a stable attractor across
a wide range of tenure values.
The vertex structure of the 129-violation coloring
mirrors the 131-violation case: $v_{42}$ remains the dominant violation hub
(co-violation degree~38, the same as in the 131-violation coloring),
confirming that the structural fault line at the extra vertex persists across tenure settings.

For comparison, the same Rust solver starting from random (Paley+FPL) seeds
achieved only 200 violations on $K_{43}$.
The bi-circulant warm-start reduced the floor by 69 violations (34\%),
demonstrating that the $\mathbb{Z}_{21}$-equivariant structure provides
meaningful guidance even for unconstrained search.

\begin{table}[h]
\centering
\begin{tabular}{lcc}
\toprule
Method & $K_{42}$ floor & $K_{43}$ floor \\
\midrule
Circulant (exhaustive)        & no valid & no valid \\
$\mathbb{Z}_{21}$ bi-circulant & 84        & 168 \\
Unconstrained tabu (random)   & 146       & 200 \\
Unconstrained tabu, tenure 12, 500 restarts & 84   & 131 \\
Unconstrained tabu, tenure 25, 200 restarts & 84   & \textbf{129} \\
Unconstrained tabu, tenure 50/100           & 84   & 129 (no improvement) \\
\bottomrule
\end{tabular}
\caption{Minimum violations under each method. Bi-circulant warm-start strictly dominates random start for both graphs; for $K_{43}$ it further escapes the bi-circulant floor. The 129-violation floor is stable across tenure values 25--100.}
\label{tab:landscape}
\end{table}

Notably, for $K_{42}$ the unconstrained search from the 84-violation bi-circulant
coloring stays at 84 (200 restarts, 39.8M iterations) — the bi-circulant floor
equals the unconstrained floor.
For $K_{43}$, the unconstrained search improves from 168 to 131, indicating
that the bi-circulant constraint is a genuine restriction in $K_{43}$:
there exist $K_{43}$ colorings with fewer violations that are not
$\mathbb{Z}_{21}$-equivariant.

\section{IRDME Structural Analysis of Ramsey Colorings}
\label{sec:irdme}

To compare the structural properties of the best colorings found, we converted
each Ramsey coloring checkpoint to a multilayer network in IRDME format
\cite{fpl2024} and analyzed them using the IRDME structural analysis pipeline.

\paragraph{Network construction.}
Each $n$-vertex coloring is represented as a three-layer network:
\begin{itemize}
\item \textbf{red\_layer}: edges colored red ($\binom{n}{2}/2$ edges approx.);
\item \textbf{blue\_layer}: edges colored blue (complement of red\_layer);
\item \textbf{violation\_layer}: pairs of vertices that co-appear in at least
      one monochromatic $K_5$ subgraph.
\end{itemize}
Each vertex receives three dimensions: \texttt{red\_degree}, \texttt{blue\_degree},
and \texttt{violation\_count} (number of monochromatic $K_5$s containing that vertex).
The FPL question becomes: do vertices that are hubs in the color layers
also dominate the violation layer?

\subsection{84-Violation $K_{42}$ Bi-Circulant}

The 84-violation $\mathbb{Z}_{21}$-equivariant coloring (Coloring~1) has
$n=42$, 1260 relations (red=420, blue=441, violation=399).
Every vertex has \texttt{red\_degree}~$= 20$ and \texttt{blue\_degree}~$= 21$ exactly
(perfect regularity from the group symmetry).
\texttt{violation\_count} ranges from 8 to 12 with mean 10.0.

IRDME law analysis: Pearson$(\text{red} \leftrightarrow \text{violation}) = 0.0000$;
all layer-pair correlations are exactly zero.
The atlas assigns every vertex rank divergence~$= 0$, placing all 42 vertices
in a single structural class. The coloring is too symmetric for IRDME to
detect any differential role.

\subsection{131- and 129-Violation $K_{43}$ Unconstrained}

The 131-violation unconstrained coloring ($n=43$, tabu tenure 12, 1518 relations:
red=453, blue=450, violation=615) breaks the perfect symmetry.
\texttt{red\_degree} ranges $[20,22]$; \texttt{violation\_count} ranges $[11,32]$
with mean 15.2, a much wider spread than the bi-circulant.
The 129-violation coloring (tabu tenure 25, same warm-start, 1529 relations:
red=453, blue=450, violation=626) has \texttt{violation\_count} ranges $[11,29]$, mean~15.0,
a slightly tighter spread as 2 violations are removed.

IRDME law: Pearson$(\text{red} \leftrightarrow \text{violation}) = +0.024$,
$p = 0.88$ (not significant). FPL is absent: color-degree hubs do not predict
violation-layer hubs.

The atlas reveals four distinct structural classes:

\begin{table}[h]
\centering
\begin{tabular}{llrc}
\toprule
Class & Vertices & Count & Characterization \\
\midrule
Universal hubs & $v_7, v_8, v_9$ & 3 & Top rank in all three layers \\
Chameleons     & $v_0, v_1, v_4, v_5, v_{10}, v_{11}, v_{20}, v_{39}, v_{42}$ & 9 & Large rank reversals across layers \\
Relay nodes    & (25 vertices) & 25 & Consistent mid-rank in all layers \\
Peripheral     & $v_{33}, v_{34}, v_{35}, v_{37}, v_{40}, v_{41}$ & 6 & Low rank in all layers \\
\bottomrule
\end{tabular}
\caption{Structural role atlas for the 131-violation $K_{43}$ coloring (tabu tenure 12). The 129-violation coloring (tenure 25) retains the same four classes with $v_{42}$ as the dominant hub.}
\label{tab:atlas}
\end{table}

The most striking result concerns $v_{42}$ (the extra vertex added to extend
$K_{42} \to K_{43}$): it is the \textbf{top hub in the violation\_layer}
(degree~38 in the co-violation graph, violation\_count~$= 32 \approx 2.1\times$ mean for 131v;
degree~38 persists in the 129-violation coloring with mean~15.0)
while being nearly peripheral in both color layers
(blue\_layer rank~\#39, red\_layer rank~\#42 of 43).
This is a direct inversion of the FPL prediction: the highest-violation
vertex has the lowest color connectivity.

The most chameleonic vertex is $v_4$: it is the \#1 hub in blue\_layer
(degree~22, maximum) yet the \#43 hub in red\_layer (degree~20, minimum) ---
a complete rank reversal between complementary layers.

\paragraph{Interpretation.}
In the symmetric bi-circulant coloring, every vertex plays an interchangeable role
and the violation structure is spread uniformly.
The unconstrained 131-violation coloring breaks this symmetry: the extra
vertex $v^*$ concentrates violations disproportionately, suggesting that the
$K_{42} \to K_{43}$ extension introduces a structural ``fault line'' that
any near-valid $K_{43}$ coloring must negotiate.
The four-class IRDME atlas is richer than any structure observable in the
symmetric coloring, exposing the heterogeneous role of vertices in the
violation landscape.

\subsection{Orbit-Type IRDME: Edge Decisions as Items}
\label{sec:irdme_orbits}

To analyze the search landscape directly, we designed a second IRDME representation
where the \emph{41 orbit types} (not the 42 vertices) are the items.
Each orbit type $T_t$ receives dimensions: current color, $K_5$ subgraph exposure
(total subgraphs containing it), active violation exposure ($K_5$s that are
currently monochromatic), and $\Delta_t$ (change in violation count if $T_t$ is flipped).
Two layers connect orbit types: \textbf{conflict\_layer} (share $\geq 1$ $K_5$ subgraph)
and \textbf{mono\_conflict\_layer} (share $\geq 1$ active violation).

Applied to the 84-violation $K_{42}$ bi-circulant:
\begin{itemize}
\item FPL law between conflict\_layer and mono\_conflict\_layer: Pearson~$= 0.000$.
      The conflict graph is nearly complete (every type shares $K_5$s with 40 others),
      leaving no variation to correlate.
\item Atlas identifies \emph{chameleon orbit types}: $T_{32}$--$T_{36}$ (orbit types
      $C_{12}$--$C_{16}$, cross-half connections) are the \#1--\#11 hubs in
      mono\_conflict\_layer (most active violations) yet ranked \#33--\#37 in
      conflict\_layer (near-peripheral in general $K_5$ overlap).
      These five cross-half orbit types are the \textbf{load-bearing decisions}
      of the 84-violation coloring: all assigned blue, all with catastrophically
      positive $\Delta_t$ (+84 to +525) if flipped to red.
\end{itemize}

\subsection{Bipartite IRDME: Vertices and Orbit Types Together}
\label{sec:irdme_bipartite}

We combined both views into a single bipartite IRDME network with
$42 + 41 = 83$ items (vertex items and orbit-type items sharing a common
\texttt{violation\_exposure} dimension) and three layers:
vertex--vertex co-violations, orbit--orbit active conflicts, and
vertex--orbit participation (vertex $v$ linked to orbit type $T_t$ if $v$
appears in a $K_5$ violation containing an edge of type $T_t$).

This design produced the \textbf{first genuine FPL signal} in any Ramsey
IRDME experiment:
\[
\text{Pearson}(\text{orbit\_conflict} \leftrightarrow \text{vertex\_orbit}) = +0.803
\quad (r^2 = 0.646, \text{ large effect}).
\]
Orbit types that are hubs in the active conflict network
also bind to more vertices through violations.

\paragraph{Cross-coloring robustness.}
We applied the bipartite IRDME to all three 84-violation $K_{42}$ bi-circulant colorings
(seeds 42, 99, and HI-42, which have different load-bearing orbit types):
the FPL signal is consistent across all three.

\begin{table}[h]
\centering
\begin{tabular}{lcc}
\toprule
Coloring & Pearson(orbit\_conflict $\leftrightarrow$ vertex\_orbit) & vertex\_orbit links \\
\midrule
Seed 42  (near-SC) & $+0.803$ & 798 \\
Seed 99  (HI, $B=A$)    & $+0.883$ & 588 \\
HI Seed 42          & $+0.768$ & 924 \\
\bottomrule
\end{tabular}
\caption{Bipartite IRDME FPL signal across the three distinct 84-violation $K_{42}$ colorings.
The orbit\_conflict $\leftrightarrow$ vertex\_orbit correlation (Pearson $+0.77$ to $+0.88$) holds
for all colorings, confirming it is a law of the 84-violation basin rather than an artifact of a
single coloring.}
\label{tab:fpl_cross}
\end{table}

The chameleon orbit types differ by coloring (seeds 42 and HI-42 share $T_{32}$--$T_{33}$ as
load-bearing; seed~99 uses $T_{37}$, $T_{39}$ instead), confirming that the 84-violation
basin is \emph{structurally broad}: multiple orbit-type configurations achieve the same
violation count via different structural compromises, yet all satisfy the same FPL law.

The complementary finding:
\[
\text{Pearson}(\text{orbit\_conflict} \leftrightarrow \text{vertex\_violation}) = -0.520,
\]
indicating that orbit conflict hubs ($C_{12}$--$C_{16}$) bind to
\emph{mid-rank vertices}, not to the vertex violation hubs ($v_0$--$v_8$).
The atlas reveals two non-overlapping violation ecosystems:
\begin{itemize}
\item \textbf{Vertex violation hubs} ($v_0$--$v_8$): top in vertex--vertex violations,
      bound to within-half orbit types $A_1, A_3, A_4, A_7, A_8$ ($T_0, T_2, T_3, T_6, T_7$).
\item \textbf{Orbit violation hubs} ($C_{12}$--$C_{16}$, $T_{32}$--$T_{36}$):
      top in orbit--orbit active conflicts, bound to mid-range vertex positions.
\end{itemize}
These two populations operate on disjoint vertex sets in the 84-violation coloring,
suggesting that the violation structure decomposes into independent subsystems.
This decomposition may constrain the space of near-valid colorings:
any improvement below 84 violations must simultaneously disrupt both subsystems.

\subsection{Topological Analysis of the Violation Complex}
\label{sec:topology}

The 129 violated $K_5$ subgraphs of the best known $K_{43}$ coloring are not independent:
each violation is a 4-simplex (five vertices, ten edges), and two violations are
\emph{adjacent} whenever they share at least one vertex.
The resulting \textbf{violation intersection graph} has 129 nodes and 4096 edges
(mean degree $63.5$, maximum $80$), forming a single connected component.

Computing the first Betti number $\beta_1 = |E| - |V| + \beta_0 = 4096 - 129 + 1 = 3968$
reveals \textbf{3968 independent cycles} in the intersection graph.
In the language of algebraic topology, the violation complex is highly non-trivial:
it cannot be contracted to a tree without cutting 3968 independent loops.

This topological measurement provides a direct explanation for the failure of
greedy descent and the stability of the 129-violation local minimum.
A single-edge flip that repairs one violated $K_5$ necessarily propagates
through these cycles, recreating violations elsewhere along connected paths.
We verified this directly: every one of the 626 ``hot'' edges (edges appearing
in at least one violation) has negative single-flip delta under the 129-violation coloring
(fixing fewer violations than it creates),
and no 2-flip combination from the five highest-coverage edges improves the count.

The dominant vertex $v_{42}$ contributes a 29-node clique to the intersection graph:
all $\binom{29}{2} = 406$ pairs of violations containing $v_{42}$ also share $v_{42}$
as a vertex, forming a complete subgraph.
Similar cliques exist for all top-participation vertices (Table~\ref{tab:atlas}).

This topological obstruction distinguishes the 129-violation basin from a
mere energy barrier: it is a structural property of the coloring that is
invariant under the choice of tabu tenure, confirmed by the observation that
tenures 25, 50, and 100 all converge to 129 violations without further improvement.
Resolving the remaining violations requires algorithms that can reason about
the global cycle structure simultaneously ---
specifically, clause-learning SAT solvers (kissat), whose learned clauses correspond
to algebraic ``cycle cuts'' that break the topological entanglement.

\paragraph{Constrained SAT formulation.}
We identified 277 ``cold'' edges (edges not participating in any violation)
and added them as unit clauses to the $K_{43}$ CNF.
Unit propagation eliminated 288 additional variables and reduced the clause count
by 80\% (from 1.93M to approximately 372k active clauses),
leaving 615 active variables.
This constrained instance encodes the problem of resolving exactly the 129 topologically
entangled violations without perturbing the 277 already-correct edges.

We ran the state-of-the-art CDCL solver kissat on this constrained instance for
2 hours 13 minutes (7994 seconds), reaching 44.3 million conflicts.
The solver found neither SAT nor UNSAT: backbone computation was abandoned after
$0.02$ seconds (no forced variables identified), 36 internal WalkSAT walks all
failed to find a satisfying assignment, and the irredundant clause database grew
from 286k to 626k without enabling further variable elimination.
This confirms that the topological obstruction (3968 independent cycles) is not
resolvable by clause learning alone within practical time limits:
the constrained neighborhood of the 129-violation coloring is neither easily
satisfiable nor refutable.

\section{Discussion and Open Questions}

The non-existence of a valid circulant coloring on $\mathbb{Z}_{42}$ suggests
that if a valid $K_{42}$ Ramsey(5,5) coloring exists, it must have
less algebraic structure than the Paley-type constructions that succeeded for
smaller Ramsey numbers. This raises the question: what is the automorphism
group of the Exoo 1989 graph?

We searched Cayley graph colorings on all three non-cyclic groups of order 42
using left-translation-orbit tabu search (20 restarts, 100k steps each,
tabu tenure 8).
The left-translation symmetry of a Cayley graph on $G$ reduces 850,668 five-subsets
to $\lfloor 850{,}668 / 42 \rfloor = 20{,}254$ canonical $K_5$ orbits ---
exactly twice the reduction achieved by the $\mathbb{Z}_{21}$ bi-circulant.
Results:

\begin{itemize}
\item $D_{21} = \mathbb{Z}_{21} \rtimes \mathbb{Z}_2$ (31 orbit variables):
      floor \textbf{84 violations} --- identical to the bi-circulant minimum.
      All violation counts are multiples of 42 (orbit size under $D_{21}$),
      with 84v $= 2 \times 42$ corresponding to exactly 2 canonical $K_5$ orbits.
\item $\mathbb{Z}_7 \times S_3$ (22 orbit variables):
      floor \textbf{84 violations} --- again the same minimum,
      reached from a smaller search space ($2^{22} \approx 4$M colorings).
\item $F_{42} = \mathbb{Z}_7 \rtimes \mathbb{Z}_6$ (Frobenius, 24 orbit variables):
      floor \textbf{210 violations} $= 5 \times 42$. The Frobenius group's
      rigid orbit structure prevents descent to the 84-violation basin
      at this tabu level.
\end{itemize}

The convergence of $D_{21}$, $\mathbb{Z}_7 \times S_3$, and $\mathbb{Z}_{21}$
bi-circulant to the same 84-violation floor strongly suggests that 84 is the
true minimum of the $K_{42}$ Ramsey violation count over all vertex-transitive
colorings, independent of the specific group symmetry.
Whether any non-vertex-transitive coloring achieves fewer than 84 violations
(or whether 84v can be driven to 0) remains open.

\paragraph{The IRDME/FPL conjecture.}
We conjecture that in any valid Ramsey(5,5;n) coloring, the distribution of
vertex degrees in the red subgraph has a specific spectral profile determined
by the FPL constraint at each vertex's neighborhood. Specifically, the vertex
degree sequence $(d_R(v))_v$ must satisfy a local Ramsey recursion:
each vertex's neighborhood is itself a valid Ramsey(4,4) coloring.
This recursive structural constraint may be exploitable as a polynomial-time
necessary condition filter for candidate colorings.

\section{Conclusion}

We have proved that no circulant graph on $\mathbb{Z}_{42}$ or $\mathbb{Z}_{43}$
provides a valid witness for Ramsey lower bounds. Both exhaustive searches used
the same bitmask $K_5$ detection algorithm (8 minutes for $K_{42}$, 7 minutes
for $K_{43}$), ruling out the most natural algebraic construction class and
establishing that any valid $K_{43}$ coloring must have non-circulant structure.

We introduced $\mathbb{Z}_{21}$-bi-circulant tabu search, which reduces the
Ramsey search problem on $K_{42}$ from $2^{861}$ binary variables to
41 orbit-type variables ($2^{41} \approx 2.2\times 10^{12}$ colorings).
Orbit-level tabu search at 180 steps/second achieves
\textbf{84 monochromatic $K_5$ violations}~--- a new computational record,
improving the previous best of 146 by 43\%.
Multiple independent seeds converge to exactly 84 violations under
$\mathbb{Z}_{21}$ bi-circulant, $\mathbb{Z}_7$ hexa-circulant, and unconstrained
tabu search symmetries. Exhaustive unconstrained search with 200 restarts
(39.8 million iterations, 26 ms/iter) starting from a non-$\mathbb{Z}_{21}$-equivariant
84-violation coloring confirmed that 84 is a stable attractor from multiple
structural basins.

We additionally characterized the 84-violation basin under symmetry constraints.
The \emph{half-identical} (HI) constraint ($B_k = A_k$) yields the same 84-violation minimum,
confirming that the 84-violation coloring can be chosen with identical within-half colorings.
The \emph{self-complementary} (SC) constraint ($B_k = 1-A_k$) reaches a floor of 126 violations
over 50,000 steps, confirming that the 84-violation basin and the SC subspace are disjoint.

We extended bi-circulant search to $K_{43}$ (43 orbit types, $2^{43}$ search space),
finding a floor of \textbf{168 violations} (Section~\ref{sec:k43bicirculant}).
The exact factor $168 = 2 \times 84$ links the $K_{42}$ and $K_{43}$ floors
under the same $\mathbb{Z}_{21}$ orbit structure.
Using the 168-violation coloring as a warm-start, unconstrained Rust tabu search
on $K_{43}$ (100 restarts, 19.8M iterations) reached \textbf{131 violations},
below the bi-circulant floor and well below the 200-violation result from random starts.
Extending to 500 restarts (248.3M total iterations, tabu tenure 12) confirmed 131 as the
floor for that tenure. Increasing the tabu tenure to 25 immediately improved the record to
\textbf{129 violations} at restart~5 (1.3M iterations),
establishing 129 as the current best known lower bound for $\min_c v(c, K_{43})$
(Table~\ref{tab:landscape}).
No search reached 0 violations; proving or disproving $R(5,5) \geq 44$ remains open.

Multilayer structural analysis (Section~\ref{sec:irdme}) reveals a
qualitative difference between the two best colorings: the symmetric
84-violation $K_{42}$ bi-circulant has all vertices in a single structural class,
while the unconstrained $K_{43}$ coloring (131 and 129 violations) exhibits four classes
and places the extra vertex $v^*$ as the top violation hub
--- a direct inversion of the FPL prediction, suggesting that the
$K_{43}$ extension creates a structural fault line that near-valid
colorings must balance.

Open directions include:
(1) Constrained SAT resolution: the 615-variable, 372k-clause instance built from
    the 129-violation cold-edge partition was run to 44.3M conflicts (2h 13min) without
    resolution, confirming that the 3968-cycle topological obstruction is not efficiently
    resolvable by CDCL; stochastic local search SAT solvers (probSAT, DRUP-based)
    or integer programming on the minimum cycle basis remain unexplored;
(2) Topological cycle decomposition: identifying a minimum cycle basis of the 3968-cycle
    violation complex and targeting those ``topological cuts'' directly with algebraic methods
    (e.g.\ integer linear programming on the hitting set formulation);
(3) Population crossover for $K_{43}$: finding a second independent 129-violation coloring
    and performing block-level crossover (analogous to the $K_{42}$ crossover experiment)
    may escape the 129-violation basin that single-trajectory tabu cannot exit;
(4) Non-vertex-transitive search on $K_{42}$: all three non-cyclic groups
    of order 42 ($D_{21}$, $\mathbb{Z}_7 \times S_3$, $F_{42}$) have been
    explored; $D_{21}$ and $\mathbb{Z}_7 \times S_3$ both achieve 84 violations,
    confirming 84 as the vertex-transitive floor. Reducing below 84 likely
    requires breaking vertex-transitivity entirely.

\paragraph{Acknowledgments.}
Computations were performed using the IRDME research platform.

\bibliographystyle{plain}
\begin{thebibliography}{9}

\bibitem{exoo1989}
G.~Exoo,
\newblock A lower bound for $R(5,5)$,
\newblock {\em J.\ Graph Theory}, 13(1):97--98, 1989.

\bibitem{angeltveit2024}
V.~Angeltveit and B.~D.~McKay,
\newblock $R(5,5) \leq 46$,
\newblock \emph{arXiv preprint arXiv:2409.15709}, 2024.

\bibitem{radziszowski}
S.~Radziszowski,
\newblock Small Ramsey numbers,
\newblock {\em Electron.\ J.\ Combin.}, Dynamic Survey DS1, updated 2021.

\bibitem{fpl2024}
V.~Sivak,
\newblock The Functional Proximity Law: hub persistence across network layers,
\newblock {\em arXiv preprint arXiv:2604.23639}, 2026.

\end{thebibliography}

\end{document}
